\subsection{Wang et al.\ (2024): efficient MTP for defects in Ni--Al alloys}
\label{sec:appendix-example-nitol2024nial}

\begin{quote}
\emph{Note on paper-id:} this corpus entry is filed under the
upstream paper-id \texttt{nitol2024nial}, but the actual authors of
the encoded paper are Wang et al.\ (Shenyang National Laboratory for
Materials Science). The CORPUS.md substitution table notes that the
original Nitol 2025 Ti-Al-V PRMaterials paper is journal-only with
no arXiv preprint and was substituted with this earlier Ni--Al MTP
work.
\end{quote}

\paragraph{Q1 -- Bibliographic identification.}
Wang, Liu, Zhu, Liu, Ma, Chen, Sun, and Chen present an efficient
moment tensor potential for defects in Ni--Al alloys, with a
genetic-algorithm-optimised tensor-contraction scheme.
arXiv:2411.01282.
\lstinputlisting[language=turtle]{artifacts/kg/listings/nitol2024nial/q01-bibliographic.ttl}

\paragraph{Q2 -- Material system.}
The MTP covers fcc Ni and L1$_2$-Ni$_3$Al, the Ni/Ni$_3$Al interface,
and their defect environments (vacancies, vacancy clusters, antisites,
GSF supercells). We model this as a single $\MaterialSystemC$ with
the defect detail in $\dlRole{materialClass}$.
\lstinputlisting[language=turtle]{artifacts/kg/listings/nitol2024nial/q02-material.ttl}

\paragraph{Q3 -- Reference calculation method.}
Reference data come from DFT calculations.
\lstinputlisting[language=turtle]{artifacts/kg/listings/nitol2024nial/q03-reference-method-type.ttl}

\paragraph{Q4 -- Reference settings.}
DFT uses VASP with PAW, PBE GGA, a 400\,eV plane-wave cutoff, and
Monkhorst--Pack k-point meshes with ${\sim}0.2$\,\AA$^{-1}$ spacing.
Convergence: total energy 1\,meV/atom, SCF $10^{-6}$\,eV, ionic-force
0.01\,eV/\AA.
\lstinputlisting[language=turtle]{artifacts/kg/listings/nitol2024nial/q04-reference-settings.ttl}

\paragraph{Q5 -- Training dataset.}
The training set has 8{,}450 configurations: 6{,}092 from on-the-fly
active-learning during AIMD plus 2{,}358 added through D-optimality
active-learning cycles. Energies, forces, and stresses are covered.
Provenance is \emph{published}.
\lstinputlisting[language=turtle]{artifacts/kg/listings/nitol2024nial/q05-dataset.ttl}

\paragraph{Q6 -- Sampling strategies.}
Three strategies were combined: on-the-fly active learning during
AIMD heating from 700 to 1600\,K (kernel-based Bayesian regression in
VASP, NPT with Langevin thermostat + Parrinello--Rahman barostat);
D-optimality active learning (configurations with extrapolation
grade $>5$ during 200\,ps MD trajectories are added; iterations
continue until the threshold is no longer exceeded); and
defect-prototype enumeration (an independent 815-configuration
validation set was constructed from the same prototypes).
\lstinputlisting[language=turtle]{artifacts/kg/listings/nitol2024nial/q06-sampling.ttl}

\paragraph{Q7 -- MLIP method, components, implementation.}
The method is a Moment Tensor Potential with a
genetic-algorithm-optimised tensor-contraction scheme: rank of moment
tensors limited to 4 ($\mu \le 3$) and a redefined level
$\text{lev}(M_{\mu,\nu}) = 2\mu + \nu + 1$ with maximum-level
threshold 2{,}653. Loss is weighted MSE on energy/forces/stresses
with weights $1{:}0.01{:}0.005$. Fitting is two-step: linear
optimisation on minimal then full basis with shared radial functions,
followed by 5{,}000 L-BFGS iterations and a least-squares
refinement. Implementation: MLIP package with custom
tensor-contraction extensions.
\lstinputlisting[language=turtle]{artifacts/kg/listings/nitol2024nial/q07-method.ttl}

\paragraph{Q8 -- Hyperparameter settings.}
Reported settings: cutoff radius $r_c=5.4$\,\AA, maximum level
2{,}653, 8 radial-basis functions, maximum body order 5, total of
2{,}653 linear basis parameters.
\lstinputlisting[language=turtle]{artifacts/kg/listings/nitol2024nial/q08-hyperparameter-settings.ttl}

\paragraph{Q9 -- MLIP run and trained model.}
A single fitting run on the 8{,}450-configuration training set
produces the optimised Ni--Al MTP. Training hardware (96 CPU cores)
is recorded on the run.
\lstinputlisting[language=turtle]{artifacts/kg/listings/nitol2024nial/q09-run-and-model.ttl}

\paragraph{Q10 -- Benchmark results.}
Validation RMSEs on the 815-configuration validation set:
1.64\,meV/atom on energy, 0.028\,eV/\AA{} on force components,
0.09\,GPa on stress components.
\lstinputlisting[language=turtle]{artifacts/kg/listings/nitol2024nial/q10-benchmark-results.ttl}

\paragraph{Q11 -- Computational resources.}
Training and inference hardware are reported as 96 CPU cores (Intel
Xeon Platinum 9242 @ 2.30\,GHz); the model achieves
${\sim}241$\,s for $10^5$ MD steps on a 2{,}048-atom Ni$_3$Al
supercell. Wall-clock training duration, peak memory, and explicit
per-atom inference time are not reported (no GPUs are used).
\lstinputlisting[language=turtle]{artifacts/kg/listings/nitol2024nial/q11-resources.ttl}

\paragraph{Q12 -- Gaps.}
Beyond Q11, the genetic-algorithm-driven tensor-contraction search
that distinguishes this MTP from the standard Shapeev formulation is
not first-class in the ontology (recorded only in the
$\dlConcept{FunctionalForm}$ label). The two-step linear-then-L-BFGS
fitting procedure is captured in $\dlRole{hasTrainingAlgorithm}$ via
prose; the ontology has no concept for multi-stage training. Many
downstream defect-property comparisons (vacancy formation energies,
GSF curves, point-defect migration barriers) are reported as derived
figures rather than encoded as separate $\BenchmarkResultC$ instances.
